%----------------------------------------------------------------------------------------
%    PACKAGES AND THEMES
%----------------------------------------------------------------------------------------

\documentclass[aspectratio=169,xcolor=dvipsnames]{beamer}

\AtBeginSection[]
{
    \begin{frame}
        \frametitle{Visão Geral}
        \tableofcontents[currentsection]
    \end{frame}
}

\usefonttheme{professionalfonts} % using non standard fonts for beamer
\usefonttheme{serif} % default family is serif
\usepackage{fontspec}
\setmainfont{XCharter}[
    UprightFont    = *-Roman,
    BoldFont       = *-Bold,
    ItalicFont     = *-Italic,
    BoldItalicFont = *-BoldItalic,
]

\usepackage{hyperref}
\usepackage{graphicx} % Allows including images
\usepackage{booktabs} % Allows the use of \toprule, \midrule and \bottomrule in tables
\input{./slides/SimplePlusTheme/beamerthemeSimplePlus.sty}
\input{./slides/SimplePlusTheme/beamerinnerthemeSimplePlus.sty}
\input{./slides/SimplePlusTheme/beamerfontthemeSimplePlus.sty}
\input{./slides/SimplePlusTheme/beamercolorthemeSimplePlus.sty}

\usepackage{listings}
\setbeamercovered{transparent}

\usepackage{biblatex}
\addbibresource{slides/reference.bib}
\renewcommand*{\bibfont}{\footnotesize}

\usepackage[brazilian ]{babel}
\usepackage{csquotes}
\usepackage{pgfplots}
\pgfplotsset{compat=1.18}

\usepackage{emoji}
\usepackage{amsmath}
\usepackage[xcharter,bigdelims,vvarbb]{newtxmath}
% newtxmath's operators font requests the fontspec family `XCharter(0)' in OT1;
% without these substitutions math digits and operator names silently fall back
% to Computer Modern (LaTeX Font Warning: Font shape `OT1/XCharter(0)/m/n' undefined).
\DeclareFontFamily{OT1}{XCharter(0)}{}
\DeclareFontShape{OT1}{XCharter(0)}{m}{n}{<->ssub * XCharter-TLF/m/n}{}
\DeclareFontShape{OT1}{XCharter(0)}{m}{it}{<->ssub * XCharter-TLF/m/it}{}
\DeclareFontShape{OT1}{XCharter(0)}{b}{n}{<->ssub * XCharter-TLF/b/n}{}
\DeclareFontShape{OT1}{XCharter(0)}{bx}{n}{<->ssub * XCharter-TLF/b/n}{}

\usepackage{bm}
\usepackage[table]{xcolor}
\usepackage{xcolor}
\usepackage{algpseudocode}

\usepackage{cancel}
\usepackage{ulem}

\usetikzlibrary{shapes}
\usetikzlibrary{arrows}
\usetikzlibrary {arrows.meta}
\usetikzlibrary{calc,positioning}
\usepgfplotslibrary{fillbetween}
\usetikzlibrary{angles,quotes}
\usetikzlibrary{tikzmark}


\definecolor{PastelRed}{HTML}{FFC1CC}
\definecolor{PastelBlue}{HTML}{C2F0FF}
\definecolor{PastelBlue2}{HTML}{3182BD}

\definecolor{PastelPink}{HTML}{FFCBE1}
\definecolor{PastelGreen}{HTML}{D6E5BD}
\definecolor{PastelYellow}{HTML}{F9E1A8}
\definecolor{PastelBlue}{HTML}{BCD8EC}
\definecolor{PastelPurple}{HTML}{DCCCEC}
\definecolor{PastelOrange}{HTML}{FFDAB4}


\renewcommand{\footnotesize}{\tiny}

\newcommand{\mespor}{
  \ifcase\month
    \or janeiro
    \or fevereiro
    \or março
    \or abril
    \or maio
    \or junho
    \or julho
    \or agosto
    \or setembro
    \or outubro
    \or novembro
    \or dezembro
  \fi
}
% \usepackage{csquotes}
% \usepackage{minted}
% \usemintedstyle{xcode}
\definecolor{vscLightBackground}{HTML}{FFFFFF}
\definecolor{vscLightText}{HTML}{000000}
\definecolor{vscLightGreen}{HTML}{008000}        % For Comments AND Numbers
\definecolor{vscLightRed}{HTML}{A31515}          % For Strings
\definecolor{vscLightTeal}{HTML}{267F99}
\definecolor{vscBlue}{HTML}{0000FF}% For main keywords (if, for, def...)
\definecolor{vscLightPytorch}{HTML}{800080}      % For PyTorch keywords (dark magenta)

\lstdefinestyle{vscode-light}{
    language=Python,
    backgroundcolor=\color{vscLightBackground},
    basicstyle=\ttfamily\scriptsize\color{vscLightText},
    keywordstyle=\color{vscLightTeal},
    commentstyle=\color{vscLightGreen},
    stringstyle=\color{vscLightRed},
    % Style for PyTorch keywords (using keyword group 2)
    keywordstyle=[2]\color{vscLightPytorch}, 
    morekeywords={self, True, False, None},
    morekeywords=[2]{
        torch, Tensor, device, dtype, to, nn, optim, autograd, utils, data,
        Module, Linear, Conv2d, ReLU, CrossEntropyLoss, Adam, SGD,
        Dataset, DataLoader, randn, zeros, ones, arange, cat, stack,
        view, reshape, sum, mean, max, min, matmul, no_grad, backward
    },
    keywordstyle=[3]\color{vscBlue}, 
    morekeywords=[3]{
        model, loss, y, y_hat, x , optimizer
    },
    frame=single,
    tabsize=1,
    showstringspaces=false,
    breaklines=true,
    captionpos=b,
    extendedchars=true,
    numbers=left,
}

\usepackage[cache=true]{minted}
% minted v3+ (TeX Live 2024+) auto-detects the output directory.
% minted v2 (Ubuntu apt TeX Live 2023) needs it set explicitly to
% match latexmkrc's $out_dir = 'build'.
\makeatletter
\@ifpackagelater{minted}{2024/01/01}{}{%
  \def\minted@outputdir{build/}%
}
\makeatother
\setminted{
    fontsize=\scriptsize,
    style=vs,
    breaklines=true,
    numbers=left,
    numbersep=5pt,
    frame=single,
}
\providecommand{\citeauthoryear}[1]{(\citeauthor{#1},~\citeyear{#1})}

\providecommand{\thetav}{\bm{\theta}}

\providecommand{\varvector}[1]{\mathbf{#1}}
% instance
\providecommand{\X}{\varvector{X}}
\providecommand{\mlinstance}{\varvector{X}}
\providecommand{\mlinstancei}{\varvector{X}^{(i)}}
\providecommand{\mlinstancex}[1]{\varvector{X}^{(#1)}}

% target
\providecommand{\y}{\varvector{y}}
\providecommand{\mltarget}{\varvector{y}}
\providecommand{\mltargeti}{\varvector{y}^{(i)}}
\providecommand{\mltargetx}[1]{\varvector{y}^{(#1)}}

%hats
\providecommand{\yhat}{\hat{\varvector{y}}}
\providecommand{\mltargethat}{\hat{\varvector{y}}}
\providecommand{\mltargethati}{\hat{\varvector{y}}^{(i)}}
\providecommand{\mltargethatx}[1]{\hat{\varvector{y}}^{(#1)}}

% linear reg
\providecommand{\mllinearregression}{\mltargethat=\mlinstance\bm{\theta}}

\providecommand{\tikzarrowcolor}[3]{%
  \begin{tikzpicture}[remember picture, overlay]
    % Arrow body
    \path[fill=#3] ([shift={(#1,#2)}]current page.south west) 
      rectangle ++(0.3, -0.2);
    % Arrow head
    \path[fill=#3] ([shift={(#1+0.3,#2+0.1)}]current page.south west) -- 
         ++(0, -0.4) -- 
         ++(0.2, 0.2) -- 
         cycle;
  \end{tikzpicture}%
}

\providecommand{\tikzarrow}[2]{%
    \tikzarrowcolor{#1}{#2}{Red}
}



\DeclareMathOperator*{\argmax}{argmax}
\DeclareMathOperator*{\argmin}{argmin}


\usetikzlibrary{fit,backgrounds}

%----------------------------------------------------------------------------------------
%    TITLE PAGE
%----------------------------------------------------------------------------------------

\title{Deep Learning}
\subtitle{Redes Recorrentes}

\author{Lucas Silveira Kupssinskü}

\institute
{
    Machine Learning Theory and Applications Laboratory \\
    PUCRS
}
\date{\mespor de \the\year}

%----------------------------------------------------------------------------------------
%    TIKZ STYLES
%----------------------------------------------------------------------------------------

\tikzset{
    rnncell/.style={draw, rounded corners, fill=PastelGreen, minimum width=1.6cm, minimum height=1.0cm},
    inputnode/.style={circle, draw, fill=PastelBlue, minimum size=0.7cm, inner sep=1pt, font=\small},
    outputnode/.style={circle, draw, fill=PastelPink, minimum size=0.7cm, inner sep=1pt, font=\small},
    opnode/.style={draw, fill=PastelYellow, rounded corners=2pt, font=\small, inner sep=2pt},
    grad/.style={->, red, very thick}
}

% Desenha uma célula recorrente com as operações internas. Os nós internos são
% posicionados em relação ao nó "+" e a caixa da célula se ajusta ao conteúdo
% (biblioteca fit), com margem igual em todos os lados. Argumentos:
%   #1 sufixo dos nomes dos nós (cell#1, xh#1, hh#1, plus#1, tanh#1, out#1,
%      xnode#1, ynode#1), pode ser vazio
%   #2 posição x do centro da célula
%   #3 rótulo da operação x·θ_ih      #4 rótulo da operação h·θ_hh
%   #5 rótulo da operação h·θ_ho      #6 rótulo do nó de entrada
%   #7 rótulo do nó de saída          #8 rótulo ŷ ao lado do nó de saída (ou vazio)
\newcommand{\rnncellgraph}[8]{%
    \node[draw, circle, inner sep=1pt, fill=white] (plus#1) at (#2, 0) {$+$};
    \node[opnode, minimum width=1.45cm, left=0.3cm of plus#1] (hh#1) {#4};
    \node[opnode, minimum width=1.45cm, right=0.3cm of plus#1] (tanh#1) {$\tanh$};
    \node[opnode, below=0.4cm of plus#1] (xh#1) {#3};
    \node[opnode, fill=PastelPink, above=0.4cm of plus#1] (out#1) {#5};
    \begin{scope}[on background layer]
        \node[rnncell, fit=(hh#1)(tanh#1)(xh#1)(out#1), inner sep=0.25cm] (cell#1) {};
    \end{scope}
    \node[inputnode] (xnode#1) at ($(plus#1 |- cell#1.south) + (0, -0.85)$) {#6};
    \node[outputnode, minimum size=0.85cm] (ynode#1) at ($(plus#1 |- cell#1.north) + (0, 0.95)$) {#7};
    \ifx\relax#8\relax\else
        \node[font=\footnotesize, anchor=south west, inner sep=1pt] at (ynode#1.north east) {#8};
    \fi
    \draw[->] (xnode#1) -- (xh#1);
    \draw[->] (xh#1) -- (plus#1);
    \draw[->] (hh#1) -- (plus#1);
    \draw[->] (plus#1) -- (tanh#1);
    \draw[->] (tanh#1) |- (out#1);
    \draw[->] (out#1) -- (ynode#1);
}

% Diagrama desdobrado de duas células usado no BPTT. As setas vermelhas do
% gradiente são desenhadas depois, em cada frame, referenciando os nós criados
% por \rnncellgraph (sufixos 1 e 2).
\newcommand{\bpttcells}{%
    \rnncellgraph{1}{0}{$1\,\theta_{ih}$}{$0\,\theta_{hh}$}{$0.462\,\theta_{ho}$}{$1$}{$0.462$}{$\hat{y}_1$}
    \rnncellgraph{2}{6.0}{$2\,\theta_{ih}$}{$0.462\,\theta_{hh}$}{$0.843\,\theta_{ho}$}{$2$}{$0.843$}{$\hat{y}_2$}
    \draw[->] (tanh1.east) -- node[above, font=\scriptsize] {$h_1\!=\!0.462$} (hh2.west);
}

%----------------------------------------------------------------------------------------
%    PRESENTATION SLIDES
%----------------------------------------------------------------------------------------

\begin{document}

\begin{frame}
    \titlepage
\end{frame}

\begin{frame}{Overview}
    \tableofcontents
\end{frame}

%========================================================================
\section{Sequências}
%========================================================================

\begin{frame}{Sequências em Aprendizado de Máquina}
\begin{itemize}
    \item Muitos problemas relevantes envolvem dados \textbf{sequenciais}, em que a ordem dos elementos importa.
    \item Exemplos:
    \begin{itemize}
        \item Análise de sentimento de uma frase ou resenha.
        \item \textit{Image captioning}: gerar uma descrição textual a partir de uma imagem.
        \item Tradução automática entre idiomas.
        \item Previsão de séries temporais (preço de ações, clima, energia).
        \item Reconhecimento de fala e síntese de áudio.
    \end{itemize}
    \item Em todos esses casos, mudar a ordem dos elementos altera o significado da entrada ou da saída.
\end{itemize}
\end{frame}

\begin{frame}{Análise de Sentimento}
\begin{itemize}
    \item A tarefa: dada uma frase, prever se o sentimento é positivo ou negativo.
\end{itemize}
\vspace{0.2em}
\begin{columns}[T]
    \begin{column}{.5\linewidth}
        \centering
        \includegraphics[height=3.0cm]{slides/03-rnn/imagens/sentimento_filme1.png}\\[0.4em]
        {\footnotesize\textit{``esse filme é uma bela porcaria''}}\;
        \raisebox{-0.35\height}{\includegraphics[height=0.55cm]{slides/03-rnn/imagens/thumb_down.png}}\\[0.3em]
        {\footnotesize sentimento \textbf{negativo}}
    \end{column}
    \begin{column}{.5\linewidth}
        \centering
        \includegraphics[height=3.0cm]{slides/03-rnn/imagens/sentimento_filme2.png}\\[0.4em]
        {\footnotesize\textit{``eu esperava que fosse uma porcaria, mas é uma bela forma de passar o tempo''}}\;
        \raisebox{-0.35\height}{\includegraphics[height=0.55cm]{slides/03-rnn/imagens/thumb_up.png}}\\[0.3em]
        {\footnotesize sentimento \textbf{positivo}}
    \end{column}
\end{columns}
\vspace{0.5em}
\begin{center}
    \colorbox{PastelOrange}{\parbox{0.9\linewidth}{\centering\small As mesmas palavras (``bela'', ``porcaria'') aparecem nas duas frases: a \textbf{ordem} faz toda a diferença!}}
\end{center}
\end{frame}

\begin{frame}{Exemplo: \textit{Image Captioning} (\textit{one-to-many})}
\begin{columns}[c]
    \begin{column}{.5\linewidth}
        \centering
        \includegraphics[width=\linewidth]{slides/03-rnn/imagens/captioning_gato.jpg}\\[0.4em]
        {\footnotesize\textit{``Um gato sentado na estrada''}}
    \end{column}
    \begin{column}{.5\linewidth}
        \centering
        \begin{tikzpicture}[every node/.style={font=\scriptsize}, scale=0.9, transform shape]
            \foreach \i in {0,1,2,3}{
                \node[rnncell, minimum width=1.0cm, minimum height=0.8cm] (c\i) at (1.7*\i, 0) {};
            }
            \foreach \i/\j in {0/1, 1/2, 2/3}{
                \draw[->, thick] (c\i) -- (c\j);
            }
            \node[inputnode, minimum size=1.0cm] (img) at (0, -1.7) {\tiny imagem};
            \draw[->, thick] (img) -- (c0);
            \foreach \i/\w in {0/Um, 1/gato, 2/sentado, 3/{na estrada}}{
                \node[outputnode] (y\i) at (1.7*\i, 1.4) {$\hat{y}_\i$};
                \draw[->, thick] (c\i) -- (y\i);
                \node[font=\scriptsize, PastelBlue2] at (1.7*\i, 2.1) {\w};
            }
        \end{tikzpicture}\\[0.5em]
        {\footnotesize Uma única entrada (a imagem) gera uma \textbf{sequência} de palavras.}
    \end{column}
\end{columns}
\end{frame}

\begin{frame}{Exemplo: Tradução Automática (\textit{many-to-many})}
\begin{columns}[c]
    \begin{column}{.54\linewidth}
        \centering
        \includegraphics[width=\linewidth]{slides/03-rnn/imagens/traducao_google.png}
    \end{column}
    \begin{column}{.46\linewidth}
        \centering
        \begin{tikzpicture}[every node/.style={font=\scriptsize}, scale=0.82, transform shape]
            \foreach \i/\inw/\outw in {0/Essa/This, 1/aula/class, 2/{é}/is, 3/{fácil}/easy}{
                \node[rnncell, minimum width=0.9cm, minimum height=0.8cm] (c\i) at (1.7*\i, 0) {};
                \node[inputnode] (x\i) at (1.7*\i, -1.3) {$x_\i$};
                \node[outputnode] (y\i) at (1.7*\i, 1.3) {$\hat{y}_\i$};
                \draw[->, thick] (x\i) -- (c\i);
                \draw[->, thick] (c\i) -- (y\i);
                \node[font=\tiny, PastelBlue2] at (1.7*\i, -1.95) {\inw};
                \node[font=\tiny, PastelBlue2] at (1.7*\i, 1.95) {\outw};
            }
            \foreach \i/\j in {0/1, 1/2, 2/3}{
                \draw[->, thick] (c\i) -- (c\j);
            }
        \end{tikzpicture}\\[0.6em]
        {\footnotesize Uma sequência de entrada gera uma sequência de saída.}
    \end{column}
\end{columns}
\end{frame}

\begin{frame}{Tipos de Problemas Sequenciais}
\centering
\begin{tikzpicture}[every node/.style={font=\tiny}, scale=0.85, transform shape]
    % one-to-one
    \node[font=\scriptsize] at (0, 1.6) {\textit{one-to-one}};
    \node[rnncell, minimum width=1.0cm, minimum height=0.7cm] (c11) at (0, 0) {};
    \node[inputnode] at (0, -1.0) {$x_t$};
    \node[outputnode] at (0, 1.0) {$\hat{y}_t$};
    \draw[->] (0,-0.6) -- (c11);
    \draw[->] (c11) -- (0, 0.6);

    % one-to-many
    \node[font=\scriptsize] at (3.5, 1.6) {\textit{one-to-many}};
    \foreach \i in {0,1,2,3}{
        \node[rnncell, minimum width=0.7cm, minimum height=0.7cm] (a\i) at (2+0.9*\i, 0) {};
    }
    \node[inputnode] at (2, -1.0) {$x_0$};
    \draw[->] (2,-0.6) -- (a0);
    \foreach \i/\j in {0/1, 1/2, 2/3}{
        \draw[->] (a\i) -- (a\j);
    }
    \foreach \i in {1,2,3}{
        \node[outputnode] at (2+0.9*\i, 1.0) {$\hat{y}_\i$};
        \draw[->] (a\i) -- (2+0.9*\i, 0.6);
    }

    % many-to-one
    \node[font=\scriptsize] at (8.0, 1.6) {\textit{many-to-one}};
    \foreach \i in {0,1,2,3}{
        \node[rnncell, minimum width=0.7cm, minimum height=0.7cm] (b\i) at (6.5+0.9*\i, 0) {};
        \node[inputnode] at (6.5+0.9*\i, -1.0) {$x_\i$};
        \draw[->] (6.5+0.9*\i,-0.6) -- (b\i);
    }
    \foreach \i/\j in {0/1, 1/2, 2/3}{
        \draw[->] (b\i) -- (b\j);
    }
    \node[outputnode] at (6.5+0.9*3, 1.0) {$\hat{y}_3$};
    \draw[->] (b3) -- (6.5+0.9*3, 0.6);

    % many-to-many
    \node[font=\scriptsize] at (12.5, 1.6) {\textit{many-to-many}};
    \foreach \i in {0,1,2,3}{
        \node[rnncell, minimum width=0.7cm, minimum height=0.7cm] (d\i) at (11+0.9*\i, 0) {};
        \node[inputnode] at (11+0.9*\i, -1.0) {$x_\i$};
        \node[outputnode] at (11+0.9*\i, 1.0) {$\hat{y}_\i$};
        \draw[->] (11+0.9*\i,-0.6) -- (d\i);
        \draw[->] (d\i) -- (11+0.9*\i, 0.6);
    }
    \foreach \i/\j in {0/1, 1/2, 2/3}{
        \draw[->] (d\i) -- (d\j);
    }
\end{tikzpicture}
\vspace{0.5em}
\begin{itemize}\footnotesize
    \item \textit{one-to-many}: \textit{image captioning} (uma imagem gera uma sequência de palavras).
    \item \textit{many-to-one}: análise de sentimento (uma sequência gera um rótulo).
    \item \textit{many-to-many}: tradução automática, geração de texto.
\end{itemize}
\end{frame}

%========================================================================
\section{Tentando Modelar com MLP e CNN}
%========================================================================

\begin{frame}{Uma Ideia: Usar um MLP}
\begin{itemize}\footnotesize
    \item Vamos usar um MLP para classificar o sentimento. Cada palavra vira uma \textit{embedding} e a camada de entrada precisa de pesos para cada palavra.
\end{itemize}
\vspace{0.2em}
\begin{columns}[T]
    \begin{column}{.5\linewidth}
        \centering
        {\footnotesize\textit{``bela porcaria''}}\\[0.4em]
        \begin{tikzpicture}[every node/.style={font=\scriptsize}, scale=0.82, transform shape]
            \node[inputnode, minimum size=0.55cm] (i1) at (0, 0.4) {bela};
            \node[inputnode, minimum size=0.55cm] (i2) at (0, -0.4) {porc.};
            \foreach \j in {1,2,3}{
                \node[draw, circle, fill=PastelGreen, minimum size=0.5cm] (h\j) at (1.9, 0.8-0.8*\j+0.8) {};
            }
            \node[draw, circle, fill=PastelGreen, minimum size=0.5cm] (o) at (3.4, 0) {};
            \foreach \i in {1,2}{ \foreach \j in {1,2,3}{ \draw[->] (i\i) -- (h\j); } }
            \foreach \j in {1,2,3}{ \draw[->] (h\j) -- (o); }
            \draw[->] (o) -- (4.0, 0);
        \end{tikzpicture}\\[0.4em]
        {\footnotesize 2 palavras: $2\times 3 = \mathbf{6}$ pesos de entrada}
    \end{column}
    \begin{column}{.5\linewidth}
        \centering
        {\footnotesize\textit{``esse filme é uma bela porcaria''}}\\[0.4em]
        \begin{tikzpicture}[every node/.style={font=\scriptsize}, scale=0.72, transform shape]
            \foreach \i/\w/\y in {1/esse/1.8, 2/filme/1.08, 3/{é}/0.36, 4/uma/-0.36, 5/bela/-1.08, 6/porc./-1.8}{
                \node[inputnode, minimum size=0.55cm, font=\tiny] (i\i) at (0, \y) {\w};
            }
            \foreach \j in {1,2,3}{
                \node[draw, circle, fill=PastelGreen, minimum size=0.5cm] (h\j) at (2.3, 0.8-0.8*\j+0.8) {};
            }
            \node[draw, circle, fill=PastelGreen, minimum size=0.5cm] (o) at (3.9, 0) {};
            \foreach \i in {1,...,6}{ \foreach \j in {1,2,3}{ \draw[->] (i\i) -- (h\j); } }
            \foreach \j in {1,2,3}{ \draw[->] (h\j) -- (o); }
            \draw[->] (o) -- (4.5, 0);
        \end{tikzpicture}\\[0.4em]
        {\footnotesize 6 palavras: $6\times 3 = \mathbf{18}$ pesos de entrada}
    \end{column}
\end{columns}
\vspace{0.3em}
\begin{center}
    \colorbox{PastelOrange}{\parbox{0.92\linewidth}{\centering\small Quanto maior a frase, mais parâmetros treináveis na \textbf{camada de entrada}, e o tamanho da entrada muda de frase para frase.}}
\end{center}
\end{frame}

\begin{frame}{Problemas com MLP em Sequências}
\begin{itemize}
    \item \textbf{Frases de tamanhos diferentes} demandariam entradas de tamanhos diferentes, com pesos diferentes.
    \item Entradas maiores demandam mais pesos.
    \item A mesma palavra em posições diferentes é multiplicada por pesos diferentes.
    \begin{itemize}
        \item Cada posição tem que aprender separadamente o significado da palavra.
    \end{itemize}
    \item Não há mecanismo natural para tratar entradas de comprimento variável.
\end{itemize}
\end{frame}

\begin{frame}{Outra Ideia: Usar uma CNN 1D}
\begin{itemize}
    \item Vamos tentar usar uma CNN para classificar o sentimento da frase.
    \item Cada palavra é convertida em uma \textit{word embedding}, e um \textit{kernel} desliza ao longo da sequência.
\end{itemize}
\vspace{0.4em}
\centering
\begin{tikzpicture}[every node/.style={font=\scriptsize}]
    % word embeddings: one row per word, 5 dimensions each
    \foreach \w [count=\i from 0] in {``o'', ``filme'', ``é'', ``muito'', ``bom''}{
        \node[anchor=east] at (-0.3, 0.8-\i*0.4) {\w};
        \foreach \j in {0,1,2,3,4}{
            \draw[fill=PastelBlue!\fpeval{30+\i*8+\j*10}!white] (\j*0.4, 0.6-\i*0.4) rectangle (\j*0.4+0.4, 1.0-\i*0.4);
        }
    }
    \node[font=\tiny] at (1.0, 1.4) {\textit{word embedding}};
    % window covered by the kernel at the first position
    \draw[dashed, thick, Gray] (-0.07, -0.27) rectangle (2.07, 1.07);

    \node[font=\normalsize] at (2.5, 0.4) {$\ast$};

    % kernel: 3 words x 5 dimensions
    \foreach \i in {0,1,2}{
        \foreach \j in {0,1,2,3,4}{
            \draw[fill=PastelYellow] (2.9+\j*0.4, 0.6-\i*0.4) rectangle (2.9+\j*0.4+0.4, 1.0-\i*0.4);
        }
    }
    \draw[thick] (2.9, -0.2) rectangle (4.9, 1.0);
    \node[font=\tiny] at (3.9, 1.4) {\textit{kernel} de tamanho 3};

    \draw[->] (5.05, 0.4) -- (5.55, 0.4);

    % activation map: one value per kernel position (5 - 3 + 1 = 3)
    \foreach \j in {0,1,2}{
        \draw[fill=red!\fpeval{40+\j*20}!white] (5.7, 0.2-\j*0.4) rectangle (6.1, 0.6-\j*0.4);
    }
    \draw[thick] (5.7, 0.2) rectangle (6.1, 0.6);
    \node[font=\tiny] at (5.9, 1.4) {\textit{activation map}};
\end{tikzpicture}
\end{frame}

\begin{frame}{Limitações da CNN 1D}
\begin{itemize}
    \item Para frases mais longas, o problema fica mais nítido:
    \begin{itemize}
        \item Contextos ``distantes'' são difíceis de modelar.
        \item Precisamos empilhar várias camadas convolucionais para aumentar o campo receptivo.
    \end{itemize}
    \item Mesmo com várias camadas, o campo receptivo é finito e o tamanho da saída depende do tamanho da entrada.
\end{itemize}
\end{frame}

\begin{frame}{Não são Boas Soluções}
\begin{itemize}
    \item MLPs exigem entradas de tamanho fixo. Frases têm tamanho variável.
    \item CNNs podem trabalhar com entradas de tamanho variável, mas se não tomarmos cuidado com o campo receptivo corremos o risco de não levar em conta a frase toda.
\end{itemize}
\vspace{0.4em}
Precisamos de uma rede que:
\begin{itemize}
    \item Trabalhe sequencialmente com as palavras, uma por vez.
    \item Mantenha um \textbf{estado} (memória) que carregue informações obtidas anteriormente.
\end{itemize}
\end{frame}

%========================================================================
\section{Recorrência em Redes Neurais}
%========================================================================

\begin{frame}{Revisitando o MLP}
\begin{columns}
    \begin{column}{.5\linewidth}
        \begin{itemize}
            \item Em uma camada densa, temos:
            \begin{itemize}
                \item entradas $x$,
                \item pesos $\theta$,
                \item função de ativação,
                \item saídas $\hat{y}$.
            \end{itemize}
        \end{itemize}
    \end{column}
    \begin{column}{.5\linewidth}
        \centering
        \begin{tikzpicture}[every node/.style={font=\small}]
            \node[inputnode] (x1) at (0, 0.7) {$x_1$};
            \node[inputnode] (x2) at (0, -0.7) {$x_2$};
            \foreach \j in {1,2,3}{
                \node[draw, circle, fill=PastelGreen, minimum size=0.65cm] (h\j) at (2.2, 1.4-\j*0.7) {};
            }
            \node[outputnode] (y) at (4.4, 0) {$\hat{y}$};
            \foreach \i in {1,2}{ \foreach \j in {1,2,3}{ \draw[->] (x\i) -- (h\j); } }
            \foreach \j in {1,2,3}{ \draw[->] (h\j) -- (y); }
            \node[font=\scriptsize, above] at (1.1, 0.95) {$\theta$};
        \end{tikzpicture}
    \end{column}
\end{columns}
\end{frame}

\begin{frame}{Revisitando o MLP}
\begin{columns}
    \begin{column}{.5\linewidth}
        \begin{itemize}
            \item Podemos \textbf{colapsar} o vetor de entrada em um único nó e a camada inteira em um único bloco.
            \item O desenho fica mais compacto, mas representa exatamente a mesma rede.
        \end{itemize}
    \end{column}
    \begin{column}{.5\linewidth}
        \centering
        \begin{tikzpicture}[every node/.style={font=\small}]
            \node[inputnode, minimum size=0.9cm] (x) at (0, 0) {$x$};
            \node[rnncell] (cell) at (2.4, 0) {};
            \node[outputnode, minimum size=0.9cm] (y) at (4.8, 0) {$\hat{y}$};
            \draw[->, thick] (x) -- (cell);
            \draw[->, thick] (cell) -- (y);
        \end{tikzpicture}
    \end{column}
\end{columns}
\end{frame}

\begin{frame}{Revisitando o MLP}
\begin{columns}
    \begin{column}{.55\linewidth}
        \begin{itemize}
            \item Agora \textbf{rotacionamos} o desenho: a entrada fica embaixo e a saída em cima.
            \item É o mesmo MLP, apenas com a representação na vertical.
            \item O vetor de entrada foi colapsado, mas continua tendo um \textbf{tamanho fixo}.
        \end{itemize}
    \end{column}
    \begin{column}{.45\linewidth}
        \centering
        \begin{tikzpicture}[every node/.style={font=\small}]
            \node[inputnode, minimum size=0.9cm] (x) at (0, -1.4) {$x$};
            \node[rnncell] (cell) at (0, 0) {};
            \node[outputnode, minimum size=0.9cm] (y) at (0, 1.4) {$\hat{y}$};
            \draw[->, thick] (x) -- (cell);
            \draw[->, thick] (cell) -- (y);
        \end{tikzpicture}
    \end{column}
\end{columns}
\end{frame}

\begin{frame}{Aplicando o MLP a uma Sequência}
\begin{itemize}
    \item Vamos rodar esse MLP na frase \textit{``Minha laranja está estragada''}, palavra a palavra.
\end{itemize}
\vspace{0.4em}
\centering
\begin{tikzpicture}[every node/.style={font=\small}]
    \foreach \i/\w in {0/Minha, 1/laranja, 2/está, 3/estragada}{
        \node[inputnode] (x\i) at (2.6*\i, -0.8) {$x_\i$};
        \node[rnncell] (c\i) at (2.6*\i, 0.4) {};
        \node[outputnode] (y\i) at (2.6*\i, 1.6) {$\hat{y}_\i$};
        \draw[->, thick] (x\i) -- (c\i);
        \draw[->, thick] (c\i) -- (y\i);
        \node[font=\scriptsize, PastelBlue2] at (2.6*\i, -1.5) {\w};
    }
\end{tikzpicture}
\vspace{0.4em}
\begin{itemize}
    \item Cada palavra é processada de forma independente.
    \item O modelo \textbf{não sabe} que ``laranja'' veio depois de ``Minha''.
\end{itemize}
\end{frame}

\begin{frame}{O Problema Continua}
\begin{itemize}
    \item Qual é o problema dessa abordagem?
    \begin{itemize}
        \item \textbf{O processamento de cada entrada é individual} e não leva em conta o restante da sequência.
        \item Não há \textbf{memória} entre os instantes.
    \end{itemize}
    \item Para um problema sequencial, precisamos que a saída em $t$ dependa também do que aconteceu em $t-1, t-2, \dots$
\end{itemize}
\end{frame}

\begin{frame}{Solução: Recorrência}
\begin{itemize}
    \item Além da entrada $\mathbf{x}_t$ no instante atual, vamos passar adiante alguma informação do instante anterior.
    \item Chamamos essa informação de \textbf{estado interno} $\mathbf{h}_t$.
\end{itemize}
\vspace{0.4em}
\centering
\begin{tikzpicture}[every node/.style={font=\small}]
    \foreach \i/\w in {0/Minha, 1/laranja, 2/está, 3/estragada}{
        \node[inputnode] (x\i) at (3.0*\i, -0.8) {$x_\i$};
        \node[rnncell] (c\i) at (3.0*\i, 0.4) {};
        \node[outputnode] (y\i) at (3.0*\i, 1.6) {$\hat{y}_\i$};
        \draw[->, thick] (x\i) -- (c\i);
        \draw[->, thick] (c\i) -- (y\i);
        \node[font=\scriptsize, PastelBlue2] at (3.0*\i, -1.5) {\w};
    }
    \foreach \i/\j in {0/1, 1/2, 2/3}{
        \draw[->, thick] (c\i) -- node[above, font=\scriptsize] {$h_\i$} (c\j);
    }
\end{tikzpicture}
\vspace{0.4em}
\begin{itemize}
    \item Os mesmos pesos são reutilizados em todos os passos de tempo.
\end{itemize}
\end{frame}

\begin{frame}{Versão Dobrada vs Desdobrada}
\begin{itemize}
    \item A representação \textbf{desdobrada no tempo} pode ser compactada em uma versão \textbf{dobrada}, com uma única célula que se conecta a si mesma.
\end{itemize}
\centering
\begin{tikzpicture}[every node/.style={font=\small}]
    % folded
    \node[inputnode] (xt) at (0, -0.8) {$x_t$};
    \node[rnncell] (ct) at (0, 0.4) {};
    \node[outputnode] (yt) at (0, 1.6) {$\hat{y}_t$};
    \draw[->, thick] (xt) -- (ct);
    \draw[->, thick] (ct) -- (yt);
    \draw[->, thick, rounded corners=4pt]
        (ct.east) -- (1.4, 0.4) -- (1.4, 2.3) -- (-1.4, 2.3) -- (-1.4, 0.4) -- (ct.west);
    \node[font=\scriptsize, right] at (1.4, 1.3) {$h_t$};

    \draw[thick] (2.3, -1.6) -- (2.3, 2.2);

    % unfolded
    \foreach \i in {0,1,2,3}{
        \node[inputnode] (ux\i) at (3.7+2.6*\i, -0.8) {$x_\i$};
        \node[rnncell] (uc\i) at (3.7+2.6*\i, 0.4) {};
        \node[outputnode] (uy\i) at (3.7+2.6*\i, 1.6) {$\hat{y}_\i$};
        \draw[->, thick] (ux\i) -- (uc\i);
        \draw[->, thick] (uc\i) -- (uy\i);
    }
    \foreach \i/\j in {0/1, 1/2, 2/3}{
        \draw[->, thick] (uc\i) -- node[above, font=\scriptsize] {$h_\i$} (uc\j);
    }
\end{tikzpicture}
\vspace{0.4em}
\begin{itemize}
    \item Os dois desenhos representam a \textbf{mesma rede}: a versão desdobrada apenas explicita os passos de tempo.
\end{itemize}
\end{frame}

\begin{frame}{Camadas Densas com Recorrência}
\begin{itemize}
    \item A saída de uma célula recorrente depende da \textbf{entrada} e do \textbf{estado interno}.
    \item Os pesos são reutilizados em todos os passos de tempo.
\end{itemize}
\vspace{0.6em}
\begin{equation*}
    \underbrace{\bigl(\,\colorbox{PastelPink}{$\hat{\mathbf{y}}_t$}\,,\;\colorbox{PastelOrange}{$\mathbf{h}_t$}\,\bigr)}_{\text{\textcolor{Magenta}{\textbf{saída}}}}
    \;=\; f\Bigl(\;
    \underbrace{\colorbox{PastelBlue}{$\mathbf{x}_t$}}_{\text{\textcolor{RoyalBlue}{\textbf{entrada}}}}\;,\;\;
    \underbrace{\colorbox{PastelOrange}{$\mathbf{h}_{t-1}$}}_{\substack{\text{\textcolor{Orange}{\textbf{memória ou}}}\\[0.1em]\text{\textcolor{Orange}{\textbf{estado interno}}}}}
    \;\Bigr)
\end{equation*}
\vspace{0.4em}
\begin{itemize}
    \item O estado interno $\mathbf{h}_{t-1}$ resume tudo o que a rede ``viu'' até o passo anterior.
\end{itemize}
\end{frame}

\begin{frame}[fragile]{Pseudocódigo de uma RNN}
\begin{lstlisting}[style=vscode-light, basicstyle=\ttfamily\scriptsize, language=Python]
rnn = RNN()
hidden_state = [0, 0, 0, 0]

sentence = ["Minha", "laranja", "está"]

for word in sentence:
    prediction, hidden_state = rnn(word, hidden_state)

# prediction -> palavra mais provavel (argmax sobre o vocabulario)
next_word = decode(prediction)
print(next_word)
# >>> "estragada"
\end{lstlisting}
\begin{itemize}
    \item A cada iteração, a rede recebe a entrada atual e o estado interno, e devolve a predição e o novo estado.
\end{itemize}
\end{frame}

\begin{frame}{Detalhando as Atualizações}
\begin{itemize}
    \item A célula recorrente mais simples (\textit{vanilla RNN}) é dada por:
\end{itemize}
\vspace{0.3em}
\begin{equation*}
    \mathbf{z}_t = \mathbf{x}_t\,\mathbf{\Theta}_{ih}^{T} + \mathbf{h}_{t-1}\,\mathbf{\Theta}_{hh}^{T} + \mathbf{b}_h, \qquad \mathbf{h}_t = \tanh(\mathbf{z}_t)
\end{equation*}
\begin{equation*}
    \hat{\mathbf{y}}_t = \mathbf{h}_t\,\mathbf{\Theta}_{ho}^{T} + \mathbf{b}_o
\end{equation*}
\begin{itemize}\footnotesize
    \item Chamamos de $\mathbf{z}_t$ a \textbf{pré-ativação} (a soma antes do $\tanh$); ela reaparece no \textit{backward pass}.
\end{itemize}
\begin{itemize}
    \item $\mathbf{\Theta}_{ih}$: pesos da entrada para o estado interno (\textit{input-to-hidden}, \texttt{weight\_ih}).
    \item $\mathbf{\Theta}_{hh}$: pesos do estado interno anterior para o novo estado interno (\textit{hidden-to-hidden}, \texttt{weight\_hh}).
    \item $\mathbf{\Theta}_{ho}$: pesos do estado interno para a saída (camada \texttt{Linear} adicional, fora do \texttt{nn.RNN}).
    \item $\mathbf{b}_h$: \textit{bias} do estado interno; \quad $\mathbf{b}_o$: \textit{bias} da saída.
\end{itemize}
\vspace{0.3em}
\begin{itemize}\footnotesize
    \item Mesma notação do \texttt{nn.RNN} do PyTorch (que usa $\mathbf{W}$ no lugar de $\mathbf{\Theta}$). O PyTorch mantém \textbf{dois} \textit{bias} no estado interno, $\mathbf{b}_{ih}$ e $\mathbf{b}_{hh}$; aqui os combinamos em um único $\mathbf{b}_h$, pois só a soma importa.
\end{itemize}
\end{frame}

%========================================================================
\section{Backpropagation Through Time}
%========================================================================

\begin{frame}{\textit{Backpropagation Through Time} (BPTT)}
\begin{itemize}
    \item O algoritmo de treinamento é o mesmo já conhecido: \textit{backpropagation}.
    \item Porém agora precisamos considerar múltiplos passos de tempo (\textit{timesteps}).
    \begin{itemize}
        \item O mesmo peso influencia a saída de várias formas diferentes, conforme o \textit{timestep}.
        \item As contribuições do gradiente são \textbf{somadas} ao longo da sequência.
    \end{itemize}
    \item Esse algoritmo é chamado de \textit{Backpropagation Through Time}.
\end{itemize}
\end{frame}

\begin{frame}{Setup do Exemplo}
\begin{columns}
    \begin{column}{.45\linewidth}
        Vamos usar a rede recorrente mais simples possível para um exemplo numérico:
        \begin{itemize}\small
            \item $x = [1, 2]$
            \item $y = 3$
            \item $h_0 = 0$
            \item $\theta_{ih} = 0.5$, \quad $\theta_{hh} = 0.5$
            \item $\theta_{ho} = 1$
            \item $b_h = 0$, \quad $b_o = 0$
        \end{itemize}
        \vspace{0.2em}
        {\footnotesize Para simplificar a aritmética, usamos dimensão 1: cada $x_t$ e $h_t$ é um escalar, e $x = [1, 2]$ são dois passos de tempo.}
    \end{column}
    \begin{column}{.55\linewidth}
        \centering
        \begin{tikzpicture}
            \rnncellgraph{}{0}{$\mathbf{x}_t\,\mathbf{\Theta}_{ih}^{T} + \mathbf{b}_h$}{$\mathbf{h}_{t-1}\,\mathbf{\Theta}_{hh}^{T}$}{$\mathbf{h}_t\,\mathbf{\Theta}_{ho}^{T} + \mathbf{b}_o$}{$\mathbf{x}_t$}{$\hat{\mathbf{y}}_t$}{}
            \draw[->] (tanh.east) -- ++(0.8, 0) node[right, font=\small] {$\mathbf{h}_t$};
        \end{tikzpicture}
    \end{column}
\end{columns}
\end{frame}

\begin{frame}{\textit{Forward Pass}, Passo 1}
\begin{itemize}\small
    \item No passo 1, $x_1 = 1$ e $h_0 = 0$:
    \begin{align*}
        h_1 &= \tanh(1\cdot\theta_{ih} + 0\cdot\theta_{hh} + b_h) = \tanh(0.5) \approx 0.462 \\
        \hat{y}_1 &= h_1\cdot\theta_{ho} + b_o = 0.462\cdot 1 + 0 = 0.462
    \end{align*}
\end{itemize}
\centering
\begin{tikzpicture}
    \rnncellgraph{}{0}{$1\,\theta_{ih} + 0$}{$0\,\theta_{hh}$}{$0.462\,\theta_{ho} + 0$}{$1$}{$0.462$}{$\hat{y}_1$}
    \draw[->] (tanh.east) -- ++(0.9, 0) node[right, font=\small] {$h_1\!=\!0.462$};
\end{tikzpicture}
\end{frame}

\begin{frame}{\textit{Forward Pass}, Passo 2}
\begin{itemize}\small
    \item No passo 2, $x_2 = 2$ e $h_1 \approx 0.462$:
    \begin{align*}
        h_2 &= \tanh(2\cdot\theta_{ih} + 0.462\cdot\theta_{hh} + b_h) = \tanh(1.231) \approx 0.843 \\
        \hat{y}_2 &= h_2\cdot\theta_{ho} + b_o = 0.843\cdot 1 + 0 = 0.843
    \end{align*}
\end{itemize}
\centering
\scalebox{0.9}{%
\begin{tikzpicture}
    \rnncellgraph{1}{0}{$1\,\theta_{ih} + 0$}{$0\,\theta_{hh}$}{$0.462\,\theta_{ho} + 0$}{$1$}{$0.462$}{$\hat{y}_1$}
    \rnncellgraph{2}{6.0}{$2\,\theta_{ih} + 0$}{$0.462\,\theta_{hh}$}{$0.843\,\theta_{ho} + 0$}{$2$}{$0.843$}{$\hat{y}_2$}
    \draw[->] (tanh1.east) -- node[above, font=\scriptsize] {$h_1\!=\!0.462$} (hh2.west);
\end{tikzpicture}}
\end{frame}

\begin{frame}{Cálculo da \textit{Loss}}
\begin{itemize}
    \item Após o \textit{forward}, comparamos $\hat{y}_2$ com a meta $y = 3$ usando metade do erro quadrático:
\end{itemize}
\begin{equation*}
    J(\theta) = \frac{1}{2}\bigl(\hat{y}_2 - y\bigr)^2 = \frac{1}{2}(0.843 - 3)^2 \approx 2.326
\end{equation*}
\begin{itemize}
    \item O fator $\tfrac{1}{2}$ apenas cancela o expoente ao derivar, semeando o \textit{backward pass}:
\end{itemize}
\begin{equation*}
    \frac{\partial J}{\partial \hat{y}_2} = (\hat{y}_2 - y)
\end{equation*}
\begin{itemize}
    \item A \textit{loss} no instante final pode ser usada sozinha (\textit{many-to-one}) ou somada ao longo de todos os instantes (\textit{many-to-many}).
\end{itemize}
\end{frame}

\begin{frame}{\textit{Backward Pass}: O Grafo Computacional}
\begin{itemize}\footnotesize
    \item O \textit{forward} desdobrado no tempo forma um \textbf{grafo computacional}: a \textit{loss} $J = \tfrac{1}{2}(\hat{y}_2 - y)^2$ depende de $\hat{y}_2$, que depende de $h_2$, que depende de $h_1$, e assim por diante.
\end{itemize}
\vspace{0.1em}
\centering
\scalebox{0.9}{%
\begin{tikzpicture}
    \bpttcells
\end{tikzpicture}}
\vspace{0.2em}
\begin{itemize}\footnotesize
    \item Para treinar, vamos percorrer esse grafo \textbf{de trás para frente}, derivando o gradiente de $J$ em relação a cada parâmetro, um de cada vez.
\end{itemize}
\end{frame}

\begin{frame}{\textit{Backward Pass}: Pesos da Saída}
\begin{columns}[T]
    \begin{column}{.5\linewidth}
        \centering
        \resizebox{\linewidth}{!}{%
        \begin{tikzpicture}
            \bpttcells
            \draw<1->[grad] (ynode2.south) to[bend left=20] (out2.north);
        \end{tikzpicture}}
    \end{column}
    \begin{column}{.5\linewidth}
        \begin{itemize}\small
            \item Começamos pelos pesos da saída: o gradiente chega em $\hat{y}_2$ e atinge a cabeça de saída (em \textcolor{red}{vermelho}).
        \end{itemize}
        \vspace{0.4em}
        {\small
        \begin{align*}
            \onslide<1->{\frac{\partial J}{\partial \theta_{ho}} &= \frac{\partial J}{\partial \hat{y}}\frac{\partial \hat{y}}{\partial \theta_{ho}} = (\hat{y} - y)\, h_2}\\[0.1em]
            \onslide<1->{&= (0.843 - 3)\cdot 0.843 \approx -1.818}\\[0.4em]
            \onslide<2->{\frac{\partial J}{\partial b_{o}} &= \frac{\partial J}{\partial \hat{y}}\frac{\partial \hat{y}}{\partial b_{o}} = (\hat{y} - y)\cdot 1}\\[0.1em]
            \onslide<2->{&= (0.843 - 3)\cdot 1 \approx -2.157}
        \end{align*}
        }
    \end{column}
\end{columns}
\end{frame}

\begin{frame}{\textit{Backward Pass}: Pesos da Entrada}
\begin{columns}[T]
    \begin{column}{.5\linewidth}
        \centering
        \resizebox{\linewidth}{!}{%
        \begin{tikzpicture}
            \bpttcells
            \draw<1->[grad] (ynode2.south) to[bend left=20] (out2.north);
            \draw<1->[grad] (out2.south) to[bend left=20] (tanh2.north);
            \draw<1->[grad] (tanh2.south) to[bend left=15] (plus2.east);
            \draw<1->[grad] (plus2.south) to[bend right=15] (xh2.north);
            \draw<2->[grad] (plus2.west) to[bend left=15] (hh2.east);
            \draw<2->[grad] (hh2.south) to[bend right=12] (tanh1.east);
            \draw<2->[grad] (tanh1.south) to[bend left=15] (plus1.east);
            \draw<2->[grad] (plus1.south) to[bend right=15] (xh1.north);
        \end{tikzpicture}}
    \end{column}
    \begin{column}{.5\linewidth}
        \begin{itemize}\small
            \item Para os pesos $\theta_{ih}$, o gradiente \textbf{atravessa todos os passos de tempo}.
        \end{itemize}
        \vspace{0.3em}
        \onslide<1->{\footnotesize%
        \[
            \frac{\partial J}{\partial \theta_{ih}} = \frac{\partial J}{\partial \hat{y}}\frac{\partial \hat{y}}{\partial h_2}\frac{\partial h_2}{\partial z_2}\!\left(\frac{\partial z_2}{\partial \theta_{ih}} + \underbrace{\frac{\partial z_2}{\partial h_1}\frac{\partial h_1}{\partial z_1}\frac{\partial z_1}{\partial \theta_{ih}}}_{\text{atravessa o tempo}}\right)
        \]}
        \onslide<2->{\footnotesize%
        \[
            \frac{\partial J}{\partial \theta_{ih}} = (\hat{y}-y)\,\theta_{ho}\,(1-\tanh^2(z_2))\bigl(x_2 + \theta_{hh}(1-\tanh^2(z_1))\,x_1\bigr)
        \]}
        \onslide<3->{\footnotesize%
        \[
            \frac{\partial J}{\partial \theta_{ih}} = (0.843 - 3)\cdot 1 \cdot (1 - 0.843^2)\bigl(2 + 0.5(1 - 0.462^2)\cdot 1\bigr) \approx -1.493
        \]}
    \end{column}
\end{columns}
\end{frame}

\begin{frame}{\textit{Backward Pass}: \textit{bias} e Pesos Recorrentes}
\begin{columns}[T]
    \begin{column}{.5\linewidth}
        \centering
        \resizebox{\linewidth}{!}{%
        \begin{tikzpicture}
            \bpttcells
            \draw<1->[grad] (ynode2.south) to[bend left=20] (out2.north);
            \draw<1->[grad] (out2.south) to[bend left=20] (tanh2.north);
            \draw<1->[grad] (tanh2.south) to[bend left=15] (plus2.east);
            \draw<2->[grad] (plus2.west) to[bend left=15] (hh2.east);
            \draw<2->[grad] (hh2.south) to[bend right=12] (tanh1.east);
            \draw<2->[grad] (tanh1.west) to[bend right=15] (hh1.east);
        \end{tikzpicture}}
    \end{column}
    \begin{column}{.5\linewidth}
        \begin{itemize}\small
            \item A mesma ideia vale para o \textit{bias} $b_h$ e para os pesos recorrentes $\theta_{hh}$.
        \end{itemize}
        \vspace{0.2em}
        {\footnotesize
        \begin{align*}
            \onslide<1->{\frac{\partial J}{\partial b_{h}} &= (\hat{y}-y)\,\theta_{ho}\,(1-\tanh^2(z_2))\bigl(1 + \theta_{hh}(1-\tanh^2(z_1))\bigr)}\\
            \onslide<1->{&\approx (0.843-3)\cdot 1 \cdot(1 - 0.843^2)\bigl(1 + 0.5(1-0.462^2)\bigr)} \\
            \onslide<1->{&\approx -0.870} \\[0.3em]
            \onslide<2->{\frac{\partial J}{\partial \theta_{hh}} &= (\hat{y}-y)\,\theta_{ho}\,(1-\tanh^2(z_2))\bigl(h_1 + \theta_{hh}(1-\tanh^2(z_1))\,h_0\bigr)} \\
            \onslide<2->{&\approx (0.843-3)\cdot 1 \cdot(1 - 0.843^2)\bigl(0.462 + 0.5(1-0.462^2)\cdot 0\bigr)} \\
            \onslide<2->{&\approx -0.288}
        \end{align*}
        }
        \onslide<2->{\begin{itemize}\footnotesize
            \item Como $h_0 = 0$, o segundo termo de $\partial J / \partial \theta_{hh}$ se anula.
        \end{itemize}}
    \end{column}
\end{columns}
\end{frame}

\begin{frame}{Resumo dos Gradientes}
\centering
\begin{tabular}{l c}
    \toprule
    Parâmetro & $\partial J / \partial \cdot$ \\
    \midrule
    $\theta_{ho}$ & $-1.818$ \\
    $b_{o}$       & $-2.157$ \\
    $\theta_{ih}$ & $-1.493$ \\
    $b_{h}$       & $-0.870$ \\
    $\theta_{hh}$ & $-0.288$ \\
    \bottomrule
\end{tabular}
\vspace{0.6em}
\begin{itemize}
    \item Repare que o gradiente do peso recorrente $\theta_{hh}$ depende dos estados internos passados.
    \item Em sequências longas, esse caminho de derivadas multiplicadas é a origem dos problemas de \textit{vanishing} e \textit{exploding gradient}.
\end{itemize}
\end{frame}

\begin{frame}{Passo do \textit{Optimizer}}
\begin{itemize}
    \item Com taxa de aprendizado $\eta = 0.1$ e os gradientes calculados:
\end{itemize}
{\small
\begin{align*}
    \theta_{ho} &\leftarrow \theta_{ho} - \eta \frac{\partial J}{\partial \theta_{ho}} = 1 - 0.1\cdot(-1.818) = 1.1818 \\
    b_{o}       &\leftarrow 0 - 0.1\cdot(-2.157) = 0.2157 \\
    \theta_{ih} &\leftarrow 0.5 - 0.1\cdot(-1.493) = 0.6493 \\
    b_{h}       &\leftarrow 0 - 0.1\cdot(-0.870) = 0.087 \\
    \theta_{hh} &\leftarrow 0.5 - 0.1\cdot(-0.288) = 0.5288
\end{align*}
}
\begin{itemize}
    \item Com os pesos atualizados, refazemos o \textit{forward} para verificar o efeito na \textit{loss}.
\end{itemize}
\end{frame}

\begin{frame}{Resultado: a \textit{Loss} Diminuiu}
\begin{itemize}\footnotesize
    \item Refazendo o \textit{forward} com os pesos atualizados ($\theta_{ih}=0.6493$, $\theta_{hh}=0.5288$, $b_h=0.087$, $\theta_{ho}=1.1818$, $b_o=0.2157$):
\end{itemize}
\centering
\scalebox{0.9}{%
\begin{tikzpicture}
    \rnncellgraph{1}{0}{$1\,\theta_{ih} + b_h$}{$0\,\theta_{hh}$}{$0.627\,\theta_{ho}$}{$1$}{$0.957$}{$\hat{y}_1$}
    \rnncellgraph{2}{6.0}{$2\,\theta_{ih} + b_h$}{$0.627\,\theta_{hh}$}{$0.938\,\theta_{ho}$}{$2$}{$1.324$}{$\hat{y}_2$}
    \draw[->] (tanh1.east) -- node[above, font=\scriptsize] {$h_1\!=\!0.627$} (hh2.west);
\end{tikzpicture}}
\vspace{0.3em}
\begin{equation*}
    J_{\text{antes}} = \tfrac{1}{2}(0.843 - 3)^2 \approx 2.326
    \;\;\longrightarrow\;\;
    J_{\text{depois}} = \tfrac{1}{2}(1.324 - 3)^2 \approx 1.404
\end{equation*}
\begin{itemize}\footnotesize
    \item Uma única iteração de \textit{BPTT} mais \textit{optimizer} já reduziu a \textit{loss}.
\end{itemize}
\end{frame}

\begin{frame}[fragile]{Implementação em PyTorch}
\begin{lstlisting}[style=vscode-light, basicstyle=\ttfamily\scriptsize, language=Python]
class RNNSimples(torch.nn.Module):
    def __init__(self) -> None:
        super().__init__()
        self.weight_ih = torch.nn.Parameter(torch.tensor(0.5))
        self.weight_hh = torch.nn.Parameter(torch.tensor(0.5))
        self.bias_h    = torch.nn.Parameter(torch.tensor(0.0))
        self.weight_ho = torch.nn.Parameter(torch.tensor(1.0))
        self.bias_o    = torch.nn.Parameter(torch.tensor(0.0))

    def forward(self, x, h=torch.tensor(0.0)):
        h = torch.tanh(x * self.weight_ih + h * self.weight_hh + self.bias_h)
        y = h * self.weight_ho + self.bias_o
        return y, h
\end{lstlisting}
\end{frame}

\begin{frame}[fragile]{Loop de Treinamento}
\begin{lstlisting}[style=vscode-light, basicstyle=\ttfamily\scriptsize, language=Python]
model = RNNSimples()
optimizer = SGD(model.parameters(), lr=0.1)
x = [torch.tensor(1.0), torch.tensor(2.0)]
y = torch.tensor(3.0)

model.train()
optimizer.zero_grad()
h = torch.tensor(0.0)
for x_i in x:
    y_hat, h = model(x_i, h)
loss = (y - y_hat)**2 / 2
loss.backward()
optimizer.step()
\end{lstlisting}
\begin{itemize}
    \item O \texttt{loss.backward()} acumula o gradiente de cada peso ao longo de todos os passos de tempo automaticamente.
\end{itemize}
\end{frame}

%========================================================================
\section{Considerações Práticas}
%========================================================================

\begin{frame}{Problemas: \textit{Vanishing} e \textit{Exploding Gradient}}
\begin{itemize}\small
    \item Os velhos conhecidos: \textit{vanishing} e \textit{exploding gradient}.
    \item Com RNNs ficam piores, pois a rede também é \textbf{profunda no tempo}.
    \begin{itemize}
        \item Em sequências longas, o gradiente se multiplica por $\theta_{hh}\,(1-\tanh^2)$ a cada passo.
        \item Após $T$ passos, o produto de Jacobianos toma a forma:
    \end{itemize}
\end{itemize}
\vspace{-0.3em}
\[
    \prod_{k=1}^{T} \theta_{hh}\,\bigl(1-\tanh^2(z_k)\bigr)
\]
\vspace{-0.6em}
\begin{itemize}\small
    \item Se cada fator é menor que 1, o produto decai exponencialmente em $T$ (\textit{vanishing}); se é maior que 1, cresce exponencialmente (\textit{exploding}).
    \item Soluções comuns:
    \begin{itemize}
        \item \textit{Gradient clipping by norm}, com limiar $c$: se $\|g\| > c$, então $g \leftarrow c \cdot g/\|g\|$. Preserva a direção do gradiente. Variantes \textit{clipping by value} distorcem a direção e são menos usadas.
        \item Arquiteturas com portas: LSTM e GRU.
    \end{itemize}
\end{itemize}
\end{frame}

\begin{frame}{Uma Observação sobre PyTorch}
\begin{columns}
    \begin{column}{.55\linewidth}
        \begin{itemize}
            \item É comum que \textit{frameworks} como o PyTorch \textbf{não} incluam a camada de saída $\hat{\mathbf{y}}$ diretamente no bloco RNN.
            \item Cabe ao desenvolvedor adicionar uma camada densa (\texttt{nn.Linear}) sobre o estado interno $\mathbf{h}_t$.
            \item Isso dá flexibilidade para combinar a célula recorrente com diferentes camadas de saída (regressão, classificação, sequência).
        \end{itemize}
    \end{column}
    \begin{column}{.45\linewidth}
        \centering
        \begin{tikzpicture}
            % mesma geometria de \rnncellgraph, sem a cabeça de saída
            \node[draw, circle, inner sep=1pt, fill=white] (plus) at (0, 0) {$+$};
            \node[opnode, minimum width=1.45cm, left=0.3cm of plus] (hh) {$\mathbf{h}_{t-1}\,\mathbf{\Theta}_{hh}^{T}$};
            \node[opnode, minimum width=1.45cm, right=0.3cm of plus] (tanh) {$\tanh$};
            \node[opnode, below=0.4cm of plus] (xh) {$\mathbf{x}_t\,\mathbf{\Theta}_{ih}^{T} + \mathbf{b}_h$};
            \begin{scope}[on background layer]
                \node[rnncell, fit=(hh)(tanh)(xh), inner sep=0.25cm] (cell) {};
            \end{scope}
            \node[inputnode] (xnode) at ($(plus |- cell.south) + (0, -0.85)$) {$\mathbf{x}_t$};
            \draw[->] (xnode) -- (xh);
            \draw[->] (xh) -- (plus);
            \draw[->] (hh) -- (plus);
            \draw[->] (plus) -- (tanh);
            \draw[->] (tanh.east) -- ++(0.8, 0) node[right, font=\small] {$\mathbf{h}_t$};
        \end{tikzpicture}
    \end{column}
\end{columns}
\end{frame}

%========================================================================
\section*{Referências e Leituras}
%========================================================================

\begin{frame}{Leituras Recomendadas}
    \leiturasize
    \begin{itemize}
        \item Capítulo 10 de \fullcite{goodfellow2016deep}.
        \item \textit{MIT Introduction to Deep Learning}, Ava Amini, 2023.
    \end{itemize}
\end{frame}

\begin{frame}[allowframebreaks]{Referências}
    \sloppy
    \printbibliography[heading=none]
\end{frame}

\end{document}
